依赖 \texttt{LLD.cpp}。对每个点，求其子树中节点数最多的相对深度；若多个深度的节点数并列最多，取最小的深度。相对深度指到该点的边数。

处理点 $x$ 时，$f[\mathrm{dfn}_x+d]$ 表示其子树中相对深度为 $d$ 的节点数。长儿子的编号为 $\mathrm{dfn}_x+1$，其相对深度 $d$ 的计数恰好位于父亲相对深度 $d+1$ 的位置，因此无需移动数组即可继承。其余儿子的计数逐项加到深度加一的位置。合并时按节点数更新答案，节点数相同时取较小的深度，最后与深度为零的当前点比较。

除长儿子外，每个儿子都是一条长链的链顶，合并时遍历的次数等于这条长链的长度。不同长链互不相交，长度之和为 $n$，所以总时间复杂度为 $O(n)$。处理父节点时会改写长儿子的部分计数，因此应在处理各点时保存答案，不能在全部处理完后再从 $f$ 中查询任意点的计数。
